Put \(r=4+\delta_0\), and condition on an arbitrary \(\mathcal F_{n,b-1}\)-history. All quantities in this paragraph are then fixed finite-population quantities. For an arm \(a\), write \(X_i=R_{nbi}(a)-\bar R_{nba}\), \(N=n_{nb}\), and \(m=m_{nba}\). By Jensen's inequality and \(\mathrm{ass:residual\text{-}moments}\),
\[
 \frac1N\sum_i|X_i|^r
 \le 2^{r-1}\left\{\frac1N\sum_i|R_{nbi}(a)|^r+|\bar R_{nba}|^r\right\}
 \le 2^r C_R. \tag{1}
\]
By \(\mathrm{ass:consistency\text{-}no\text{-}interference}\), the observed residual on a label assigned to arm \(a\) equals \(R_{nbi}(a)\). By \(\mathrm{ass:complete\text{-}block\text{-}randomization}\), those labels form a simple random sample without replacement of size \(m\) from the \(N\) labels.

We first derive the finite-population transport bound used below. Let
\[
 \nu=\frac1m\sum_{i:Z_{nbi}=a}\delta_{X_i},
 \qquad
 \mu=\frac1N\sum_i\delta_{X_i}.
\]
Clip both measures to \([-T,T]\), partition that interval into \(H\) cells of length at most \(2T/H\), and project each cell to one representative. If \(P_h\) and \(\widehat P_h\) are the population and sample masses of cell \(h\), the exact first- and second-order inclusion probabilities of simple random sampling give
\[
 \operatorname{Var}(\widehat P_h\mid\mathcal F_{n,b-1})
 =\frac{N-m}{m(N-1)}P_h(1-P_h)\le\frac1{2m},
\]
where \(N\ge2\). Thus Cauchy--Schwarz implies
\[
 E\!\left(\sum_{h=1}^H|\widehat P_h-P_h|\mid\mathcal F_{n,b-1}\right)
 \le \frac{H}{\sqrt{2m}}. \tag{2}
\]
Couple mass that lies in the same cell and move unmatched clipped mass by at most \(2T\). Equations (1)--(2), together with \(X_i^2\mathbf 1\{|X_i|>T\}\le |X_i|^rT^{-(r-2)}\), yield
\[
 E\{W_2^2(\nu,\mu)\mid\mathcal F_{n,b-1}\}
 \le C\left(\frac{T^2}{H^2}+\frac{T^2H}{\sqrt m}+T^{-(r-2)}\right). \tag{3}
\]
Indeed, the population clipping cost is at most \(2^rC_RT^{-(r-2)}\), and the conditional expected sample clipping cost equals the population clipping cost by the first-order inclusion probability. Choose \(H=\lceil m^{1/6}\rceil\) and \(T=m^{1/(3r)}\). Every term in (3) is then bounded by a common constant times
\[
 m^{-\beta},\qquad \beta=\frac{r-2}{3r}. \tag{4}
\]
The plug-in margin translates \(\nu\) by minus its sample mean. Since \(N^{-1}\sum_iX_i=0\), the simple-random-sampling covariance identity gives
\[
 E\left[\left\{m^{-1}\sum_{i:Z_{nbi}=a}X_i\right\}^2\middle|\mathcal F_{n,b-1}\right]
 =\left(\frac1m-\frac1N\right)S_{nba}^2\le C/m.
\]
Translation by a scalar has squared quadratic-transport cost equal to the squared scalar. The triangle inequality for \(W_2\), (3)--(4), and the preceding display therefore imply
\[
 E\{W_2^2(\widehat\mu_{nba},\mu_{nba})\mid\mathcal F_{n,b-1}\}
 \le C m_{nba}^{-\beta}. \tag{5}
\]

For the variance term, write \(\overline{X^2}_S=m^{-1}\sum_{i:Z_{nbi}=a}X_i^2\) and \(\overline{X^2}=N^{-1}\sum_iX_i^2\). The same inclusion-probability covariance identity, now applied to \((X_i^2)_i\), and equation (1) give
\[
 E\{|\overline{X^2}_S-\overline{X^2}|\mid\mathcal F_{n,b-1}\}
 \le \left\{\left(\frac1m-\frac1N\right)
 \frac1{N-1}\sum_i(X_i^2-\overline{X^2})^2\right\}^{1/2}
 \le Cm^{-1/2}.
\]
Moreover, the preceding sample-mean identity gives \(E(\bar X_S^2\mid\mathcal F_{n,b-1})\le C/m\). Since
\[
 \widehat S_{nba}^2=\frac{m}{m-1}(\overline{X^2}_S-\bar X_S^2),
 \qquad S_{nba}^2=\frac{N}{N-1}\overline{X^2},
\]
the triangle inequality and \(|m/(m-1)-1|\le2/m\), \(|N/(N-1)-1|\le2/N\) for \(m,N\ge2\) give
\[
 E\{|\widehat S_{nba}^2-S_{nba}^2|\mid\mathcal F_{n,b-1}\}
 \le C m_{nba}^{-1/2}. \tag{6}
\]
All constants in (1)--(6) depend only on \(\delta_0\) and \(C_R\).

We next derive continuity of both multi-marginal endpoints. For each arm, take an optimal quadratic coupling of \(\mu_{nba}\) and \(\widehat\mu_{nba}\). Because these measures have finite support, disintegrating those finitely supported couplings and drawing the second coordinates conditionally independently glues them to any \(\gamma\in\Gamma_{nb}\). If \(X\) has law \(\gamma\) and \(Y\) is the resulting vector with empirical margins, then \(\mathrm{ass:contrast\text{-}normalization}\) and Cauchy--Schwarz give
\[
 \begin{aligned}
 |E(c^\top X)^2-E(c^\top Y)^2|
 &\le\{E(c^\top(X-Y))^2\}^{1/2}\{E(c^\top(X+Y))^2\}^{1/2}\\
 &\le\left\{\sum_aW_2^2(\mu_{nba},\widehat\mu_{nba})\right\}^{1/2}
       \left\{2\sum_a(E X_a^2+E Y_a^2)\right\}^{1/2}.
 \end{aligned}\tag{7}
\]
Apply (7) first to an optimizer for the population margins and then, with the roles reversed, to an optimizer for the empirical margins. Such optimizers exist because the relevant finite-dimensional coupling polytopes are nonempty and compact. Taking conditional expectations, applying Cauchy--Schwarz, using (5), and using the common second-moment bound following from (1), we obtain for both signs
\[
 E\{|\widehat J_{nb}^{\pm}-J_{nb}^{\pm}|\mid\mathcal F_{n,b-1}\}
 \le C\left(\min_a m_{nba}\right)^{-\beta/2}. \tag{8}
\]
This applies to minima and maxima and does not require a unique transport plan.

Let \(t_{nb}=w_{nb}^2/n_{nb}\), so \(\sum_bt_{nb}=s_n^2\). Executed overlap gives \(m_{nba}^{-1}\le(\varepsilon_0n_{nb})^{-1}\), while \((n_{nb}-1)^{-1}\le2n_{nb}^{-1}\). Combining (6)--(8) with the definitions of both endpoints, summing absolute block errors, and applying the tower property along the original adaptive filtration gives, for all sufficiently large \(n\),
\[
 \begin{aligned}
 \sup_{\mathbb D\in\mathfrak D}
 E_{\mathbb D}\frac{|\widehat L_n-L_n|+|\widehat U_n-U_n|}{s_n^2}
 &\le \frac{C}{s_n^2}\sum_bt_{nb}
 \left(q_n^{-1/2}+q_n^{-\beta/2}\right)\\
 &\le Cq_n^{-\kappa},
 \qquad \kappa=\frac{\beta}{2}=\frac{2+\delta_0}{6(4+\delta_0)}.
 \end{aligned}\tag{9}
\]
No union bound over blocks or histories appears in (9). At block \(b\), the predictable counts and fitted values are held fixed only at the pre-block history, the conditional bound is uniform at that history, and the tower property then integrates that history out. In particular, neither \(\log(2B_nK)\) nor the leverage envelope is used.

For all sufficiently large \(n\), (9) is stronger than the claimed truncated-expectation bound. For the remaining finitely many \(n\), enlarge the common constant: the outer truncation is at most one and \(q_n^{-\kappa}>0\). Thus
\[
 \sup_{\mathbb D\in\mathfrak D}E_{\mathbb D}\left[1\wedge
 \frac{|\widehat L_n-L_n|+|\widehat U_n-U_n|}{s_n^2}\right]
 \le C_{\mathrm{ep}}q_n^{-\kappa}=C_{\mathrm{ep}}r_n.
\]
All constants are uniform over \(\mathfrak D\) and depend only on \(K,\varepsilon_0,\delta_0,C_R\). If \(X_n\) denotes the nonnegative normalized endpoint error, then
\(\mathbf 1\{X_n>\eta\}\le(1\wedge X_n)/(1\wedge\eta)\); taking expectations proves the displayed probability bound. Finally, the Hausdorff distance between two real intervals is \(\max\{|\widehat L_n-L_n|,|\widehat U_n-U_n|\}\), which is bounded by their summed endpoint error. Since \(q_n\to\infty\) and \(\kappa>0\), \(r_n=q_n^{-\kappa}\to0\). Every probability and expectation above is under the original sequential randomization law, with no conditioning on \(\mathcal F_{n,B_n}\).